% Page 003 — Table des matières
% Page liminaire : volontairement NON numérotée.
%
% La pagination du document commencera avec la page 005,
% qui porte le numéro imprimé 1.


\begin{center}
  {\Large\bfseries MEYRUEIS AU FIL DU TEMPS}
\end{center}

\vspace{1.0cm}

\noindent
\begin{tabularx}{\textwidth}{|P{2.7cm}|Y|P{1.5cm}|}
\hline

\textbf{1622}
&
\textbf{Avant-propos}\newline
Lettre d'Henri Belon, colporteur\newline
Carte du pays de Meyrueis
&
1--2\newline
3--4
\\ \hline

\textbf{1700-1705}
&
\textbf{Trois figures marquantes :} Jacques Crégut le camisard,\newline
Roman le prédicant et Bertrand de Campis le galérien
&
5--14
\\ \hline

\textbf{1768}
&
\textbf{Épidémie à Meyrueis}
&
15--17
\\ \hline

\textbf{1789-1793}
&
\textbf{Les foires de Meyrueis}
&
18
\\ \hline

\textbf{1794}
&
\textbf{La Société Populaire de Meyrueis}
&
19
\\ \hline

\textbf{1789-1795}
&
\textbf{Protestants et Catholiques face à la Révolution}\newline
Histoire du pasteur-administrateur Samuel François
&
20--23
\\ \hline

\textbf{1884-1894}
&
\textbf{Histoire de Meyrueis, La fabrique de pointes, « Nous\newline
avons remonté la vallée de schiste bleu... »}
&
24--27
\\ \hline

\textbf{1900}
&
\textbf{Les chapeliers de Meyrueis}
&
28--31
\\ \hline

\textbf{1920}
&
\textbf{La grotte de la Douce}
&
32--36
\\ \hline

\textbf{1919-1939}
&
\textbf{Meyrueis entre les deux guerres}
&
37--38
\\ \hline

\textbf{1933-1935}
&
Souvenirs du pasteur Jean Gounelle.
&
39--46
\\ \hline

\textbf{1940-1945}
&
\textbf{Les Chantiers de Jeunesse, Georges Roth, Albert Girah}
&
47--59
\\ \hline

\textbf{1940-1944}
&
Notes Marceline Causse : Meyrueis 1942--1944
&
60--61
\\ \hline

\textbf{1985}
&
\textbf{Anecdotes et témoignage Éva Verdier}
&
62--63
\\ \hline

&
\textbf{La Vallée de la Jonte}
&
64
\\ \hline

&
\textbf{Le château d'Ayres}
&
65--69
\\ \hline

&
\textbf{Le moulin d'Ayres} \hfill \textbf{Gatuzières}
&
70--75
\\ \hline

&
\textbf{Cabrillac}
&
76
\\ \hline

&
\textbf{La Vallée de la Brèze}
&
77
\\ \hline

&
\textbf{L'arbre du poivre}
&
78
\\ \hline

&
\textbf{La Société Cévenole de Carburants forestiers}
&
79--86
\\ \hline

\textbf{1943}
&
\textbf{Cabanals, Alauze, Raffègues, Pourcarès Colliergues...}
&
87--92
\\ \hline

\textbf{1935}
&
\textbf{Les Fêtes des Oubrets}
&
93--95
\\ \hline

\textbf{1938}
&
\textbf{Instituteur aux Oubrets}
&
96--110
\\ \hline

\textbf{1930-1964}
&
\textbf{La Vallée du Béthuzon Ferrussac et Les Harkis Le Villaret}
&
110--121
\\ \hline

&
\textbf{Meyrueisiens connus ou à connaître}
&
122--126
\\ \hline

&
\textbf{Epilogue}
&
127--128
\\ \hline

\end{tabularx}

